%!TEX root = main.tex
%!TEX spellcheck = en-US

We start by playing ecumenical games over the Kripke semantics we have just defined. 

Embedding Kripke semantics into intuitionistic game models bridges the gap between semantic evaluation (evaluating truth based on possible worlds) and interactive proof theory (evaluating truth via a dialogue or game between two players). In a classical Hintikka game, players evaluate a formula in a single, static model. In contrast, to evaluate an intuitionistic formula game-theoretically, the game must take place over a Kripke frame $(W, \le)$, where worlds represent states of information or knowledge that accumulate over time.

Extending Hintikka dialogue games~\cite{Hintikka1973-HINLLA-2} to logics supported by Kripke semantics is conceptually straightforward~\cite{blackburn-synthese,Robert-PhD}: in addition to the current roles of the players and the current formula $F$, one only has to  keep track of the current world $w$ in the model. However, this extension introduces an unfortunate drawback: the {\em game trees}, \ie, labelled trees whose nodes are game states, are no longer determined solely by the syntax of the formula, but instead depend on the relational structure of the model. This is in stark contrast to semantic games for propositional classical logic, where semantic information is required only at the final stage to determine the winner. The loss of uniformity in game trees across all models is a significant limitation of this approach.

As in~\cite{blackburn-synthese,DBLP:conf/wollic/Freiman21,DBLP:conf/csl/PimentelOLFF25}, we address this problem by allowing explicit references to worlds and the accessibility relation within the object language. Let $\A=\{\ag,\b,\ldots\}$ be a non-empty set of agents,
$\At=\{p,q,\ldots\}$ be a countable set of propositional variables, and $N=\{i,j,\ldots\}$ be a countable set of \emph{nominals}. The language of NAME~is generated by the following grammar
$$F  ::= p  \mid \neg F  \mid F_1  \wedge F_2  \mid F_1  \vee F_2  \mid  R(i,j) $$
where $p\in \At$, and $i,j\in N$.  Formulas of the form $p$, $R(i,j)$ are called \emph{elementary}. We shall use $F,G,H$ to range over formulas. 
The propositional connectives $\top$, $\bot$, $\to$ can be obtained in the usual way. 

Intuitively, nominals are used as names for worlds of the model, while the propositions $R^\pm(i,j)$ state that agent $i$ is a  \emph{friend}/\emph{enemy} (or, more generally, \emph{agrees}/\emph{disagrees}) with
 $j$. 


A model $\mathbb{M}$ is a tuple $\langle \A,\R,\V,\g\rangle$ where $\A$
is a set (of agents), $\g:N\rightarrow \A$ is called \emph{denotation
function}, $\R\subseteq \A\times \A$, and $\V:\At\rightarrow
\mathcal{P}(\A)$. 
%
A model is a NAME-model if $\g$ is surjective, \ie, every agent has a name.  

A formula $F$ is true over $\mathbb{M}$, written $\mathbb{M}\Vdash F$ iff
$\mathbb{M},\ag\Vdash F$, for all agent $\ag\in\A$. For a set of formulas
$\Delta$, we write $\mathbb{M}\models \Delta$ iff $\mathbb{M}\Vdash \Delta$ for all $F
\in \Delta$. A formula $F$ is valid iff $\mathbb{M}\Vdash F$ for every
NAME-model $\mathbb{M}$. For a class of models $\mathfrak{M}$, we write
$\Delta\models_{\mathfrak{M}} F$ iff $\mathbb{M}\Vdash F$ for every model
$\mathbb{M}\in \mathfrak{M}$ with $\mathbb{M}\models \Delta$. 
